% ID: SEL-D0006-OA-E01
% ESTADO: validada. S007, 2026-11-07; tres pasadas del mismo autor y revisión matemática, didáctica y visual.
\textbf{Variables y restricciones.}\par\nopagebreak[4]
Sean $x$ e $y$ los comprimidos diarios de UNAL y DOSAL. El coste es $K=0.6x+y$ euros. Cada mínimo nutricional impone una desigualdad con signo mayor o igual:
\[
5x+2y\geq10,\qquad5x+5y\geq20,\qquad x+3y\geq6,\qquad x,y\geq0.
\]
La segunda restricción equivale a $x+y\geq4$. La región factible es cerrada y no acotada, por encima de las tres rectas y dentro del primer cuadrante.
\begin{center}\includegraphics[width=0.55\linewidth]{figuras/s007-region.pdf}\end{center}
Las intersecciones factibles consecutivas de la frontera inferior son
\[
(0,5),\quad(2/3,10/3),\quad(3,1),\quad(6,0).
\]
La primera y última son cortes con ejes. Las dos interiores se obtienen de $5x+2y=10$ con $x+y=4$, y de $x+y=4$ con $x+3y=6$, respectivamente. La intersección de calcio y calorías es $(18/13,20/13)$: se descarta porque $x+y=38/13<4$.

\textbf{Coste y mínimo global.}\par\nopagebreak[4]
En esos vértices, $K$ vale $5$, $56/15$, $14/5$ y $18/5$ euros. El candidato óptimo es $(3,1)$. Para probarlo en toda la región, incluidas sus partes no acotadas, combinamos dos restricciones:
\[
K=\frac25(x+y)+\frac15(x+3y)\geq\frac25(4)+\frac15(6)=\frac{14}{5}.
\]
La igualdad obliga a $x+y=4$ y $x+3y=6$, de donde $x=3$, $y=1$. Por tanto,
\[
\boxed{3\text{ comprimidos UNAL y }1\text{ DOSAL};\quad K_{\min}=2.8\text{ euros diarios}.}
\]
Comprobación nutricional: aporta $5(3)+2=17$ unidades de calcio, $5(3)+5=20$ de proteínas y $3+3=6$ calorías. Cumple los tres mínimos. Aunque se estudió primero el modelo continuo, el óptimo tiene coordenadas enteras y también resuelve el problema con comprimidos completos. El sombreado sólo representa una ventana de la región, que continúa hacia arriba y derecha.
